% Part: first-order-logic
% Chapter: natural-deduction
% Section: proving-things-quant

\documentclass[../../../include/open-logic-section]{subfiles}

\begin{document}

\olfileid[fr]{fol}{ntd}{prq}

\olsection{\usetoken{P}{derivation} avec quantificateurs}

\begin{ex}
Avec les quantificateurs, il faut veiller à respecter la condition
sur l'eigenvariable. Cela oblige parfois à modifier l'ordre de
certaines inférences. En général, il est utile de préparer d'abord
les règles soumises à cette condition : elles se trouveront plus
bas dans la preuve achevée.

Construisons une !!{derivation} de la !!{formula}
$\lexists[x][\lnot !A(x)] \lif \lnot \lforall[x][!A(x)]$.
Comme d'habitude, nous commençons par écrire :
\begin{prooftree}
\AxiomC{}
\UnaryInfC{$\lexists[x][\lnot !A(x)]\lif \lnot \lforall[x][!A(x)]$}
\end{prooftree}
Précisons d'abord ce qu'il faudrait pour justifier la dernière
étape par la règle \Intro{\lif}.
\begin{prooftree}
\AxiomC{$\Discharge{\lexists[x][\lnot !A(x)]}{1}$}
\DeduceC{$\lnot \lforall[x][!A(x)]$}
\DischargeRule{\Intro{\lif}}{1}
\UnaryInfC{$\lexists[x][\lnot !A(x)]\lif \lnot \lforall[x][!A(x)]$}
\end{prooftree}
Aucune règle à appliquer à $\lnot \lforall[x][!A(x)]$ ne
s'impose immédiatement; préparons donc la !!{derivation} de façon
à utiliser \Elim{\lexists}. Il faut ici respecter la condition
sur l'eigenvariable : choisissons une constante absente de
$\lexists[x][\lnot !A(x)]$, de la conclusion
$\lnot \lforall[x][!A(x)]$ et des autres hypothèses dont dépend
l'inférence. Si aucune constante ne figure dans les formules
concernées, toute constante convient; sinon, nous en choisissons
une qui n'y figure pas.
\begin{prooftree}
\AxiomC{$\Discharge{\lexists[x][\lnot !A(x)]}{1}$}
\AxiomC{$\Discharge{\lnot !A(a)}{2}$}
\DeduceC{$\lnot \lforall[x][!A(x)]$}
\DischargeRule{\Elim{\lexists}}{2}
\BinaryInfC{$\lnot \lforall[x][!A(x)]$}
\DischargeRule{\Intro{\lif}}{1}
\UnaryInfC{$\lexists[x][\lnot !A(x)] \lif \lnot \lforall[x][!A(x)]$}
\end{prooftree}
Pour dériver $\lnot \lforall[x][!A(x)]$, essayons la règle
\Intro{\lnot}. Il faut dériver une contradiction, en utilisant
éventuellement $\lforall[x][!A(x)]$ comme hypothèse supplémentaire.
Cette contradiction peut bien sûr dépendre de l'hypothèse
$\lnot !A(a)$, qui sera déchargée par l'inférence
\Elim{\lexists}. On peut organiser la construction ainsi :
\begin{prooftree}
\AxiomC{$\Discharge{\lexists[x][\lnot !A(x)]}{1}$}
\AxiomC{$\Discharge{\lnot !A(a)}{2}, \Discharge{\lforall[x][!A(x)]}{3}$}
\DeduceC{$\lfalse$}
\DischargeRule{\Intro{\lnot}}{3}
\UnaryInfC{$\lnot \lforall[x][!A(x)]$}
\DischargeRule{\Elim{\lexists}}{2}
\BinaryInfC{$\lnot \lforall[x][!A(x)]$}
\DischargeRule{\Intro{\lif}}{1}
\UnaryInfC{$\lexists[x][\lnot !A(x)]\lif \lnot \lforall[x][!A(x)]$}
\end{prooftree}
Nous sommes près d'obtenir une contradiction. La règle la plus
simple à appliquer est \Elim{\lforall}, qui n'impose aucune
condition sur une eigenvariable. Nous pouvons remplacer la
variable universellement quantifiée~$x$ par n'importe quel terme
clos. Il est donc judicieux d'utiliser encore~$a$ pour obtenir
la contradiction recherchée.
\begin{prooftree}
\AxiomC{$\Discharge{\lexists[x][\lnot !A(x)]}{1}$}
\AxiomC{$\Discharge{\lnot !A(a)}{2}$}
\AxiomC{$\Discharge{\lforall[x][!A(x)]}{3}$}
\RightLabel{\Elim{\lforall}}
\UnaryInfC{$!A(a)$}
\RightLabel{\Elim{\lnot}}
\BinaryInfC{$\lfalse$}
\DischargeRule{\Intro{\lnot}}{3}
\UnaryInfC{$\lnot \lforall[x][!A(x)]$}
\DischargeRule{\Elim{\lexists}}{2}
\BinaryInfC{$\lnot \lforall[x][!A(x)]$}
\DischargeRule{\Intro{\lif}}{1}
\UnaryInfC{$\lexists[x][\lnot !A(x)]\lif \lnot \lforall[x][!A(x)]$}
\end{prooftree}
Il faut maintenant vérifier, surtout en présence de quantificateurs,
qu'aucune condition sur une eigenvariable n'a été enfreinte. La
seule règle appliquée qui impose une telle condition est
\Elim{\lexists}. La constante~$a$ choisie n'apparaît ni dans sa
prémisse existentielle, ni dans sa conclusion, ni dans les autres
hypothèses concernées. La !!{derivation} est donc correcte.
\end{ex}

\begin{ex}
Il arrive que nous dérivions une !!{formula} à partir d'autres
!!{formula}s. Certaines hypothèses peuvent alors rester non
déchargées. Il faut donc suivre les hypothèses utilisées aussi
soigneusement que la conclusion recherchée.

Construisons une !!{derivation} de la !!{formula}
$\lexists[x][!C(x,b)]$ à partir des hypothèses
$\lexists[x][(!A(x) \land !B(x))]$ et
$\lforall[x][(!B(x) \lif !C(x,b))]$.
Comme d'habitude, écrivons la conclusion tout en bas.
\begin{prooftree}
\AxiomC{}
\UnaryInfC{$\lexists[x][!C(x,b)]$}
\end{prooftree}
Nous disposons de deux hypothèses. Pour exploiter la première,
c'est-à-dire chercher une !!{derivation} de
$\lexists[x][!C(x, b)]$ en exploitant l'hypothèse
$\lexists[x][(!A(x) \land !B(x))]$, avec l'aide de la seconde
hypothèse, nous utilisons
\Elim{\lexists}. Cette règle étant soumise à une condition sur
l'eigenvariable, préparons d'abord son application. Nous choisissons
une constante~$a$ distincte de~$b$ et absente des hypothèses ainsi
que de la conclusion. Nous obtenons :
\begin{prooftree}
\AxiomC{$\lexists[x][(!A(x) \land !B(x))]$}
\AxiomC{$\Discharge{!A(a) \land !B(a)}{1}$}
\DeduceC{$\lexists[x][!C(x,b)]$}
\DischargeRule{\Elim{\lexists}}{1}
\BinaryInfC{$\lexists[x][!C(x,b)]$}
\end{prooftree}
Les deux hypothèses de départ font intervenir~$!B$. Il peut être
utile d'appliquer maintenant \Elim{\land} pour obtenir~$!B(a)$.
\begin{prooftree}
\AxiomC{$\lexists[x][(!A(x) \land !B(x))]$}
\AxiomC{$\Discharge{!A(a) 
\land !B(a)}{1}$}
\RightLabel{\Elim{\land}}
\UnaryInfC{$!B(a)$}
\DeduceC{$\lexists[x][!C(x,b)]$}
\DischargeRule{\Elim{\lexists}}{1}
\BinaryInfC{$\lexists[x][!C(x,b)]$}
\end{prooftree}
La seconde hypothèse dont nous disposons est
$\lforall[x][(!B(x) \lif !C(x,b))]$. La règle \Elim{\lforall}
ne comporte pas de condition sur une eigenvariable : nous pouvons
instancier $x$ par la !!{constant}~$a$ et obtenir
$!B(a) \lif !C(a, b)$. Nous avons alors à la fois
$!B(a) \lif !C(a,b)$ et $!B(a)$. L'étape suivante est une
application directe de \Elim{\lif}.
\begin{prooftree}
\AxiomC{$\lexists[x][(!A(x) \land !B(x))]$}
\AxiomC{$\lforall[x][(!B(x) \lif !C(x,b))]$}
\RightLabel{\Elim{\lforall}}
\UnaryInfC{$!B(a) \lif !C(a,b)$}
\AxiomC{$\Discharge{!A(a) 
\land !B(a)}{1}$}
\RightLabel{\Elim{\land}}
\UnaryInfC{$!B(a)$}
\RightLabel{\Elim{\lif}}
\BinaryInfC{$!C(a,b)$}
\DeduceC{$\lexists[x][!C(x,b)]$}
\DischargeRule{\Elim{\lexists}}{1}
\insertBetweenHyps{\hspace{-5em}}
\BinaryInfC{$\lexists[x][!C(x,b)]$}
\end{prooftree}
Il ne reste qu'une application de \Intro{\lexists} pour atteindre
notre but.
\begin{prooftree}
\AxiomC{$\lexists[x][(!A(x) \land !B(x))]$}
\AxiomC{$\lforall[x][(!B(x) \lif !C(x,b))]$}
\RightLabel{\Elim{\lforall}}
\UnaryInfC{$!B(a) \lif !C(a,b)$}
\AxiomC{$\Discharge{!A(a) 
\land !B(a)}{1}$}
\RightLabel{\Elim{\land}}
\UnaryInfC{$!B(a)$}
\RightLabel{\Elim{\lif}}
\BinaryInfC{$!C(a,b)$}
\RightLabel{\Intro{\lexists}}
\UnaryInfC{$\lexists[x][!C(x,b)]$}
\DischargeRule{\Elim{\lexists}}{1}
\insertBetweenHyps{\hspace{-5em}}
\BinaryInfC{$\lexists[x][!C(x,b)]$}
\end{prooftree}
Nous avons veillé à chaque étape à respecter les conditions sur
les eigenvariables. La !!{derivation} obtenue est donc correcte.
\end{ex}

\begin{ex}
Construisons une !!{derivation} de la !!{formula}
$\lnot\lforall[x][!A(x)]$ à partir des hypothèses
$\lforall[x][!A(x)] \lif \lexists[y][!B(y)]$ et
$\lnot\lexists[y][!B(y)]$.
Comme d'habitude, écrivons la !!{formula} recherchée tout en bas.
\begin{prooftree}
\AxiomC{}
\UnaryInfC{$\lnot\lforall[x][!A(x)]$}
\end{prooftree}
La dernière ligne de la !!{derivation} est une négation : essayons
donc \Intro{\lnot}. Il faudra trouver comment !!{derive} une
contradiction.
\begin{prooftree}
\AxiomC{$\Discharge{\lforall[x][!A(x)]}{1}$}
\DeduceC{$\lfalse$}
\DischargeRule{\Intro{\lnot}}{1}
\UnaryInfC{$\lnot\lforall[x][!A(x)]$}
\end{prooftree}
La construction est bien engagée. Nous pourrions appliquer
\Elim{\lforall}, mais il n'est pas évident que cela nous rapproche
du but. Utilisons plutôt l'une de nos hypothèses :
$\lforall[x][!A(x)] \lif \lexists[y][!B(y)]$, jointe à
$\lforall[x][!A(x)]$, permet d'appliquer \Elim{\lif}.
\begin{prooftree}
\AxiomC{$\lforall[x][!A(x)] \lif \lexists[y][!B(y)]$}
\AxiomC{$\Discharge{\lforall[x][!A(x)]}{1}$}
\RightLabel{\Elim{\lif}}
\BinaryInfC{$\lexists[y][!B(y)]$}
\DeduceC{$\lfalse$}
\DischargeRule{\Intro{\lnot}}{1}
\UnaryInfC{$\lnot\lforall[x][!A(x)]$}
\end{prooftree}
Il nous reste une hypothèse à exploiter; elle permet d'obtenir
une contradiction par \Elim{\lnot}.
\begin{prooftree}
\AxiomC{$\lnot\lexists[y][!B(y)]$}
\AxiomC{$\lforall[x][!A(x)] \lif \lexists[y][!B(y)]$}
\AxiomC{$\Discharge{\lforall[x][!A(x)]}{1}$}
\RightLabel{\Elim{\lif}}
\BinaryInfC{$\lexists[y][!B(y)]$}
\RightLabel{\Elim{\lnot}}
\BinaryInfC{$\lfalse$}
\DischargeRule{\Intro{\lnot}}{1}
\UnaryInfC{$\lnot\lforall[x][!A(x)]$}
\end{prooftree}
\end{ex}

\begin{prob}
Donner des !!{derivation}s établissant les résultats suivants :
\begin{enumerate}
\item $\Proves (\lforall[x][!A(x)] \land \lforall[y][!B(y)]) \lif
\lforall[z][(!A(z) \land !B(z))]$.
\item $\Proves (\lexists[x][!A(x)] \lor \lexists[y][!B(y)]) \lif
\lexists[z][(!A(z) \lor !B(z))]$.
\item $\lforall[x][(!A(x) \lif !B)] \Proves \lexists[y][!A(y)] \lif !B$.
\item $\lforall[x][\lnot !A(x)] \Proves \lnot\lexists[x][!A(x)]$.
\item $\Proves \lnot\lexists[x][!A(x)] \lif \lforall[x][\lnot !A(x)]$.
\item $\Proves \lnot\lexists[x][\lforall[y][((!A(x,y) \lif \lnot
!A(y,y)) \land (\lnot !A(y,y) \lif !A(x,y)))]]$.
\end{enumerate}
\end{prob}

\begin{prob}
Donner des !!{derivation}s établissant les résultats suivants :
\begin{enumerate}
\item $\Proves \lnot\lforall[x][!A(x)] \lif \lexists[x][\lnot!A(x)]$.
\item $(\lforall[x][!A(x)] \lif !B) \Proves \lexists[y][(!A(y) \lif !B)]$.
\item $\Proves \lexists[x][(!A(x) \lif \lforall[y][!A(y)])]$.
\end{enumerate}
(Tous ces résultats nécessitent la règle~$\FalseCl$.)
\end{prob}

\begin{editorial}
Dans le premier exemple, la source suppose qu'aucune constante ne
figure dans les formules; pour une formule schématique arbitraire,
il faut plutôt choisir une constante fraîche. La prémisse
existentielle est écrite ici avec sa négation, conformément à
l'arbre, et le rappel de la condition inclut la conclusion.
Tout terme utilisé pour l'élimination universelle reste clos.
Dans le deuxième exemple, la fraîcheur de $a$, notamment par
rapport à $b$, est explicitée; les deux hypothèses de départ
restent disponibles pendant la construction. Une parenthèse mal placée dans
la prémisse existentielle du troisième arbre de cet exemple est
rétablie avant le crochet fermant l'argument du quantificateur.
\end{editorial}

\end{document}
